%%% 经济学讲义元信息 —— 改书名、作者、版本只需改这一处
\title{经济学原理讲义}
\subject{Principles of Economics}
\author{Clarence Liao}
\institute{}
\htversion{0.1.0}
\date{\today}
\license{本讲义以 CC BY-NC-SA 4.0 协议发布，转载请注明出处。}
\shorttitle{经济学原理}

%%% 封面图：生产可能性边界。取自第 1 章图 1.1，选它是因为它一张图说完
%%% 本书的出发点——资源有限，边界之内是浪费，边界之外够不着，
%%% 一切选择都发生在那条弧上，取此舍彼。
\coverart{%
  \begin{tikzpicture}[scale=0.62]
    %% 可行域
    \fill[htRule, opacity=0.07] (0,0) -- (5,0) arc[start angle=0, end angle=90, radius=5] -- cycle;
    %% 坐标轴
    \draw[htGray, -{Stealth[length=4.5pt]}] (0,0) -- (6.3,0);
    \draw[htGray, -{Stealth[length=4.5pt]}] (0,0) -- (0,6.3);
    %% 生产可能性边界
    \draw[htRule, line width=1.3pt] (5,0) arc[start angle=0, end angle=90, radius=5];
    %% 三个点：边界上、内部、外部
    \fill[htAccent] (3,4) circle (0.085);
    \fill[htGray]   (2,2) circle (0.085);
    \fill[htAccent] (4.5,3.5) circle (0.085);
    \node[htAccent, font=\scriptsize, anchor=south west] at (3.05,4.05) {\(A\)};
    \node[htGray,   font=\scriptsize, anchor=north east] at (1.95,1.95) {\(B\)};
    \node[htAccent, font=\scriptsize, anchor=south west] at (4.55,3.55) {\(C\)};
  \end{tikzpicture}}
